\section{直线和圆的位置关系}

\subsection{直线和圆的联立}

\begin{theorem}[直线与圆的交点的弦长]
    然后利用圆心到直线的距离 $d$、半径 $r$ 和弦长 $L$ 之间的关系 $L=2\sqrt{r^2-d^2}$。
\end{theorem}

\begin{example}
    已知直线 $l:(m-1) x+2 y+3-m=0$ 与圆 $C: x^2+y^2-6 x+6 y=0$ 交于 $A, B$ 两点，则线段 $A B$ 的长度的取值范围是（ \hspace{2em} ）
    \begin{tasks}(4)
        \task $[\sqrt{10}, 3 \sqrt{2}]$
        \task $[2 \sqrt{10}, 6 \sqrt{2}]$
        \task $[2 \sqrt{10}, 4 \sqrt{2}]$
        \task $[\sqrt{10}, 6 \sqrt{2}]$
    \end{tasks}
\end{example}

\vspace{4em}

\subsection{弦长公式}

\begin{theorem}[弦长公式]
    设直线和圆的两个交点为 $A(x_1, y_1), B(x_2, y_2)$，圆心为 $O$，则弦长 $L$ 满足
    \[
        L = \sqrt{(x_2 - x_1)^2 + (y_2 - y_1)^2}
    \]
    设这条直线的斜率是 $k$，则：

\end{theorem}

\begin{theorem}[直线与圆的交点及韦达定理]
    将直线方程代入圆的方程，得到关于交点某一坐标（如 $x$）的一元二次方程。利用韦达定理求出两交点横坐标之积 $x_M x_N$，再结合向量点积的坐标表示 $\overrightarrow{O M} \cdot \overrightarrow{O N} = (1+k^2)x_M x_N$ 求解。
\end{theorem}
\begin{example}
    直线 $y=k x$ 与圆 $(x-1)^2+(y-1)^2=1$ 交于 $M 、 N$ 两点，$O$ 为坐标原点，则 $\overrightarrow{O M} \cdot \overrightarrow{O N}=$（ \hspace{2em} ）
    \begin{tasks}(4)
        \task $\frac{1}{1+k^2}$
        \task $\frac{k^2}{1+k^2}$
        \task $1$
        \task $2$
    \end{tasks}
\end{example}

\vspace{4em}

\subsection{直线与圆的距离问题}

\begin{theorem}[点到直线的距离与圆的几何性质]
    求圆上动点与固定两点构成三角形面积的最值问题，通常转化为求圆上动点到该固定两点所在直线的距离的最值问题。其最小（大）距离等于圆心到直线的距离减（加）去半径。
\end{theorem}
\begin{example}
    已知两点 $A(-1,0), ~ B(0,2)$ ，点 $P$ 是圆 $(x-2)^2+y^2=1$ 上任意一点，则 $\triangle P A B$ 面积的最小值是（ \hspace{2em} ）
    \begin{tasks}(4)
        \task $\sqrt{5}$
        \task $2$
        \task $3+\frac{\sqrt{5}}{2}$
        \task $3-\frac{\sqrt{5}}{2}$
    \end{tasks}
\end{example}

